\documentclass[a4paper,12pt]{article}

\usepackage{fontspec}
\setmainfont{Times New Roman}
\usepackage[russian]{babel}
\usepackage{graphicx, float, url}
\usepackage{geometry}
\geometry{left=3cm,right=1.5cm,top=2cm,bottom=2cm}
\usepackage{setspace}
\onehalfspacing
\setlength{\parindent}{1.25cm}

\begin{document}

\begin{titlepage}
  \thispagestyle{empty}
  \begin{center}

    
\includegraphics[width=0.18\textwidth]{../assets/gerb.png}

    \vspace{0.5cm}

    {\small
      МИНИСТЕРСТВО НАУКИ И ВЫСШЕГО ОБРАЗОВАНИЯ\\
      РОССИЙСКОЙ ФЕДЕРАЦИИ
    }

    \vspace{0.3cm}

    {\small
      Федеральное государственное бюджетное образовательное учреждение\\
      высшего образования
    }

    \vspace{0.3cm}

    {\large \textbf{«МИРЭА --- Российский технологический университет»}}

    {\normalsize РТУ МИРЭА}

    {\normalsize Институт искусственного интеллекта (ИИИ)}\\
    {\normalsize Кафедра промышленной информатики (ПИ)}

    \vspace{2cm}

    {\LARGE \textbf{ОТЧЁТ}}

    \vspace{0.3cm}

    {\Large \textbf{ПО ПРАКТИЧЕСКОЙ РАБОТЕ № 20}}

    \vspace{0.4cm}

    {\large по дисциплине «Информатика»}

    \vspace{2.5cm}

    \begin{tabular}{ll}
      Выполнил студент группы КВБО-11-25 & Шапаренко Ф.А. \\
      Принял ассистент: & Беднов Г.А. \\
    \end{tabular}

    \vfill

    Практическую работу выполнил: «23» апреля 2026 г.\\
    «Зачтено» \rule{1cm}{0.4pt} 2026 г.

    \vspace{1cm}

    Москва 2026

  \end{center}
\end{titlepage}

\noindent\textbf{Тема:} ПЕРВИЧНЫЙ И РАЗВЕДОЧНЫЙ АНАЛИЗ ДАННЫХ\\
\noindent\textbf{Цель работы:} \\ 1. Ознакомится с концепцией
первичного и разведочного анализа данных.  \\ 2. Построить различные
виды диаграмм данных по варианту. Сделать  заключение о характере
предоставленных данных.  \\ 3. Рассчитать описательные характеристики
данных по варианту.

Вариант 27
\begin{figure}[H]
  \centering
  \includegraphics[width=1\textwidth]{/home/shafed/projects/02-computer-science/defend-4/assets/Figure_1.png}
\end{figure}

\begin{figure}[H]
  \centering
  \includegraphics[width=1\textwidth]{/home/shafed/projects/02-computer-science/defend-4/assets/Figure_2.png}
\end{figure}

\begin{figure}[H]
  \centering
  \includegraphics[width=1\textwidth]{/home/shafed/projects/02-computer-science/defend-4/assets/Figure_3.png}
\end{figure}

\begin{figure}[H]
  \centering
  \includegraphics[width=1\textwidth]{/home/shafed/projects/02-computer-science/defend-4/assets/Figure_4.png}
\end{figure}

\begin{figure}[H]
  \centering
  \includegraphics[width=1\textwidth]{/home/shafed/projects/02-computer-science/defend-4/assets/Figure_5.png}
\end{figure}

\paragraph{Заключение.}
Чем меньше цена, тем больше продаж — обратная линейная зависимость. \\
Диапазон цен: от 624,95 до 1 511,25 \\
Диапазон продаж: от 125 до 169 \\
Большинство сделок совершается по ценам от 807 до 936 руб. при объёме
от 152 до 160 единиц.

\end{document}
